\documentclass[10pt]{article}
\usepackage[margin=1in]{geometry}
\usepackage{graphicx}
\graphicspath{{./images/}}

\usepackage{amsmath, amssymb, amsthm, amsfonts, mathtools, bbm, breqn}
\usepackage{enumitem}       % for shortlabels
\usepackage{multicol}       % multiple columns
\usepackage{abraces}        % asymmetric braces
\usepackage{skull}
\usepackage{tikz}
\usetikzlibrary{arrows.meta, arrows, calc, matrix, positioning}

\usepackage{hyperref}
\usepackage[capitalize]{cleveref}


% Citing theorems by name. (source: https://tex.stackexchange.com/questions/109843/cleveref-and-named-theorems)
\makeatletter
\newcommand{\ncref}[1]{\cref{#1}\mynameref{#1}{\csname r@#1\endcsname}}

\def\mynameref#1#2{%
  \begingroup
  \edef\@mytxt{#2}%
  \edef\@mytst{\expandafter\@thirdoffive\@mytxt}%
  \ifx\@mytst\empty\else
  \space(\nameref{#1})\fi
  \endgroup
}
\makeatother

% for the pipe symbol
\usepackage[T1]{fontenc}

\usepackage{xcolor} % Enables a broader range of colors

% Define custom colors
\definecolor{RoyalBlue}{cmyk}{1, 0.50, 0, 0}

% Note commands
\definecolor{Red}{rgb}{1,0,0}
\definecolor{Blue}{rgb}{0,0,1}
\definecolor{Purple}{rgb}{.75,0,.25}
\newcommand{\rnote}[1]{\textcolor{Red}{[#1]}}       % Red note
\newcommand{\pnote}[1]{\textcolor{Purple}{[#1]}}    % Purple note
\newcommand{\bnote}[1]{\textcolor{Blue}{#1}}        % Blue text

\usepackage{xparse} % allows paragraph-spaning color environment
\NewDocumentCommand{\Max}{+m}{%
  {\color{Red}#1}%
}

% Emphasized text
\newcommand{\demph}[1]{\textcolor{RoyalBlue}{\textbf{\slshape #1}}} % Slanted


% Claim numbering (the counter restarts after each proof environment)
\newcounter{claimcount}
\setcounter{claimcount}{0}
\newenvironment{claim}{\refstepcounter{claimcount}\par\addvspace{\medskipamount}\noindent\textbf{Claim \arabic{claimcount}:}}{}
\usepackage{etoolbox}
\AtBeginEnvironment{proof}{\setcounter{claimcount}{0}}
\newenvironment{claimproof}{\par\addvspace{\medskipamount}\noindent\textit{Proof of Claim  \arabic{claimcount}.}}{\hfill\ensuremath{\qedsymbol} \tiny{Claim}

  \medskip}
% Add claim support to cleverref
\crefname{claimcount}{Claim}{Claims}


% Math Environments
\newtheorem{theorem}{Theorem}
\newtheorem{assumption}[theorem]{Assumption}
\newtheorem{lemma}[theorem]{Lemma}
\newtheorem{proposition}[theorem]{Proposition}
\newtheorem{corollary}[theorem]{Corollary}
\newtheorem{question}[theorem]{Question}
\theoremstyle{definition}
\newtheorem{definition}[theorem]{Definition}
\newtheorem{remark}[theorem]{Remark}
\newtheorem{example}[theorem]{Example}
\newtheorem{notation}[theorem]{Notation}
\newtheorem{problem}[theorem]{Problem}

% Redefine the Example environment to include "End of example [number]"
\makeatletter
\let\oldexample\example
\renewenvironment{example}
{\begin{oldexample}}
  {\par\smallskip\hfill   End of Example~\theexample. $\square$    \par\end{oldexample}}
\makeatother

% Custom Math Commands
\newcommand{\vt}{\vspace{5mm}}                       % Vertical space
\newcommand{\fl}[1]{\noindent\textbf{#1}}            % Bold first line
\newcommand{\Fl}[1]{\vspace{5mm}\noindent\textbf{#1}}% Bold first line with space above
\newcommand{\norm}[1]{\left\lVert #1 \right\rVert}   % Norm

% Common math symbols
\newcommand{\R}{\mathbb{R}}           % Real numbers
\newcommand{\N}{\mathbb{N}}           % Natural numbers
\newcommand{\C}{\mathbb{C}}           % Complex numbers
\newcommand{\Z}{\mathbb{Z}}           % Integers
\newcommand{\Q}{\mathbb{Q}}           % Rational numbers
\newcommand{\E}{\mathbb{E}}           % Expectation
\renewcommand{\P}{\mathbb{P}}         % Probability (renamed to avoid \P clash)
\newcommand{\indep}{\perp\!\!\!\perp} % Independence symbol

% Operators
\newcommand{\Var}{\mathrm{Var}}   % Variance
\newcommand{\tr}{\mathrm{tr}}     % Trace
\newcommand{\dist}{\mathrm{dist}} % Distance


\usepackage{biblatex}
\addbibresource{refs.bib}


\begin{document}
\begin{center}
  \subsection*{Quiz 1 Study Guide}
\end{center}
\textit{\textbf{The quiz will be in-class on Friday, March 13. It will be 15
    minutes long.} Some combination of the following problems will be on the
  quiz. You do not need to turn this study guide in; it is purely for your
  benefit. On the exam, please aim to answer the questions as precisely as you
  can (when grading your answers, I will be looking for mathematical
  definitions, not examples).} \vspace{1em}

\section{Questions from before quiz 1 (some of these will be on quiz 2)}

\begin{problem}
  What does it mean for two events $A$ and $B$ to be \demph{mutually exclusive}?
\end{problem}

\begin{problem}
  In order for a function $\P$ to be a \demph{probability measure}, it must
  satisfy three axioms. State all three axioms.
\end{problem}


\begin{problem}
  Let $A$ and $B$ be events, and assume that $\P\left[A \right]>0$. What is the
  \demph{conditional probability of $B$ given $A$}?
\end{problem}

\begin{problem}
  What is a \demph{$k$-tuple}?
\end{problem}

\begin{problem}
  Describe the difference between a \demph{combination} and a \demph{permutation}.
\end{problem}

\begin{problem}
  A \demph{set} is an \underline{\phantom{\quad \quad \quad \quad}} collection
  of elements, all of which are \underline{\phantom{\quad \quad \quad \quad}}.
\end{problem}

\begin{problem}
  What does it mean for two events to be \demph{disjoint}?
\end{problem}

\begin{problem}
  What is the \demph{sample space} of an experiment? What is an \demph{event}?
\end{problem}


\begin{problem}
  What does it mean for a sequence of events $E_{1},E_{2},E_{3},\ldots $ to be
  \demph{pairwise disjoint}?
\end{problem}

\begin{problem}
  Let $A$ and $B$ be events. Define the following terms
  \begin{itemize}
    \item the \demph{complement} of $A$
    \item the \demph{union} of $A$ and $B$
    \item the \demph{intersection} of $A$ and $B$
  \end{itemize}
\end{problem}

\begin{problem}
  What does it mean for to events $A$ and $B$ to have a \demph{positive relationship}?
  A \demph{negative relationship}?
\end{problem}


\section{New stuff since quiz 1}

\begin{problem}
  What does it mean for $n$ events $A_{1},\ldots,A_{n}$ to be \demph{mutually exhaustive}?
\end{problem}

\begin{problem}
  What does it mean for two events to be \demph{independent}? What does it
  mean for them to be \demph{dependent}?
\end{problem}

\begin{problem}
  What is a \demph{random variable}?
\end{problem}

\begin{problem}
  What does it mean for $n$ events $A_{1},\ldots,A_{n}$ to be \demph{mutually independent}?
\end{problem}





\begin{problem}
  What does it mean for a random variable to be \demph{discrete}?
\end{problem}

\begin{problem}
  What is the \demph{probability mass function} (pmf) of a discrete random variable?
\end{problem}

\begin{problem}
  What is the \demph{cumulative distribution function} (cdf) of a random variable $X$?
\end{problem}



\begin{problem}
  Suppose $X$ is a discrete random variable. What is the \demph{expected
    value} of $X$?
\end{problem}


\begin{problem}
  What is \demph{linearity of expectation}?
\end{problem}


\begin{problem}
  What is the \demph{variance} of a random variable? The \demph{standard deviation}?
\end{problem}


\begin{problem}
  What is a \demph{binomial} random variable? Make sure to state the probability
  mass function.
\end{problem}


\begin{problem}
  What is a \demph{geometric} random variable? Make sure to state the
  probability mass function.
\end{problem}


\begin{problem}
  What is a \demph{Poisson} random variable? Make sure to state the
  probability mass function.
\end{problem}

\begin{problem}
  What is a \demph{probability density function}?
\end{problem}

\begin{problem}
  What does it mean for a random variable to be \demph{continuous}?
\end{problem}

\begin{problem}
  What is a \demph{uniform} random variable? Make sure to write down the
  pdf.
\end{problem}

\begin{problem}
  What is an \demph{exponential} random variable? Make sure to write down the
  pdf.
\end{problem}

\begin{problem}
  What is the \demph{expected value} of a continous random variable?
\end{problem}
\end{document}